\documentclass[11pt]{article}

\usepackage[margin=1in]{geometry}
\usepackage{amsmath,amssymb,amsthm}
\usepackage{mathtools}
\usepackage{stmaryrd}      % \llbracket, \rrbracket
\usepackage{mathpartir}    % \inferrule for SOS rules
\usepackage{listings}
\usepackage{xcolor}
\usepackage{booktabs}
\usepackage{enumitem}
\usepackage[hidelinks]{hyperref}
\usepackage{comment}
\usepackage{tikz}
\usetikzlibrary{fit,positioning,arrows.meta}

%% ---- notation ----
\newcommand{\Loc}{\mathit{Loc}}
\newcommand{\Var}{\mathit{Var}}
\newcommand{\Val}{\mathit{Val}}
\newcommand{\Tid}{\mathit{Tid}}
\newcommand{\ThId}{\mathit{ThId}}   % thread identifiers (the scheduled thread i)
\newcommand{\TrId}{\mathit{TxId}}   % transaction identifiers (the live transaction t)
\newcommand{\Cmd}{\mathit{Cmd}}
\newcommand{\Expr}{\mathit{Expr}}
\newcommand{\store}{\sigma}
\renewcommand{\log}{\lambda}
\newcommand{\reduces}{\longrightarrow}
\newcommand{\tred}{\mathrel{\longrightarrow_{t}}}
\newcommand{\unit}{()}
\newcommand{\bbegin}{\textsf{begin}}
\newcommand{\eend}{\textsf{end}}
\newcommand{\retry}{\textsf{retry}}
\newcommand{\fork}{\textsf{fork}}
\newcommand{\getc}[1]{\textsf{get}\;(#1)}
\newcommand{\setc}[2]{\textsf{set}\,(#1,#2)}
\newcommand{\nin}{\bot}                 % "no transaction" marker
\newcommand{\subst}[3]{#1[#2\mapsto #3]}
\newcommand{\rulename}[1]{\textsc{\small #1}}

%% ---- evaluation-context notation ----

\newcommand{\ectx}{\mathcal{H}}
\newcommand{\ehole}[1]{\ectx[[\, #1 \,]]}
%% runtime get/set actions: they record the value read/written and the
%% transaction token t, written get(l,v)@t and set(l,v)@t.
\newcommand{\getev}[3]{\textsf{get}(#1,#2)\,@\,#3}
\newcommand{\setev}[3]{\textsf{set}(#1,#2)\,@\,#3}
\newcommand{\elided}{\mathrm{actions}}
\newtheorem{definition}{Definition}
\theoremstyle{remark}
\newtheorem*{remark}{Remark}
\theoremstyle{plain}

%% ---- abbreviations used in the worked-examples section ----
\newcommand{\st}[1]{[\,l\,{\mapsto}\,#1\,]}          % store over the single location l
\newcommand{\emap}[1]{\{#1\}}                          % thread map E
\newcommand{\adm}{\xrightarrow{\;\mathrm{adm}\;}}     % administrative reduction step
\newcommand{\Abody}{\textsf{let}\;x\,{=}\,\getc{l}\;\textsf{in}\;\setc{l}{x\,{+}\,1}}
\newcommand{\Bbody}{\textsf{let}\;y\,{=}\,\getc{l}\;\textsf{in}\;
   (\textsf{if}\;y\,\texttt{==}\,0\;\textsf{then}\;(\setc{l}{99}\,;\,\retry)\;
   \textsf{else}\;\setc{l}{y\,{+}\,1})}

%% ---- declarative-semantics notation (Cerone et al. style) ----
\newcommand{\Op}{\mathit{Op}}
\newcommand{\EventId}{\mathit{EventId}}
\newcommand{\HEvent}{\mathit{HEvent}}
\newcommand{\WEvent}{\mathit{WEvent}}
\newcommand{\REvent}{\mathit{REvent}}
\newcommand{\rdev}[2]{\mathsf{read}(#1,#2)}
\newcommand{\wrev}[2]{\mathsf{write}(#1,#2)}
\newcommand{\retryev}{\rho}                 % retry event marker
\newcommand{\po}{\mathit{po}}
\newcommand{\maxpo}{\mathrm{max}_{\po}}
\newcommand{\att}{\mathit{att}}
\newcommand{\anyop}[2]{\mathord{\_}(#1,#2)} % read(x,n) or write(x,n)
\newcommand{\pob}{\xrightarrow{\po}}        % program-order edge

%% ---- listings style for the Reagent language ----
\definecolor{kw}{rgb}{0.10,0.10,0.55}
\definecolor{cm}{rgb}{0.30,0.45,0.30}
\lstdefinelanguage{reagent}{
  keywords={let,in,get,set,begin,end,retry,if,then,else},
  sensitive=true,
  morecomment=[l]{//},
  morecomment=[n]{(*}{*)},
  morestring=[b]",
}
\lstset{
  language=reagent,
  basicstyle=\ttfamily\small,
  keywordstyle=\color{kw}\bfseries,
  commentstyle=\color{cm}\itshape,
  showstringspaces=false,
  columns=fullflexible,
  frame=single,
  framesep=4pt,
  xleftmargin=4pt,
}
\newcommand{\code}[1]{\lstinline[basicstyle=\ttfamily]{#1}}

\title{\textbf{The Reagent Language}\\[2pt]}
\date{\today}

\begin{document}
\maketitle

\begin{abstract}
TODO
\end{abstract}

\section{Abstract syntax}

\[
\begin{array}{rcll}
n     & \in & \mathbb{N}        & \text{natural-number literals}\\
l     & \in & \Loc              & \text{shared mutable locations}\\
x     & \in & \Var              & \text{immutable, let-bound variables}\\[4pt]
v     & ::= & n \mid \unit \mid \textsf{true} \mid \textsf{false} & \text{values (naturals, unit, booleans)}\\[4pt]
c     & ::= & \getc{l}          & \text{read a shared location}\\
      & \mid & \setc{l}{e}      & \text{write a shared location}\\
      & \mid & \bbegin          & \text{open a transaction}\\
      & \mid & \eend            & \text{commit the current transaction}\\
      & \mid & \retry           & \text{abort and restart the transaction}\\[4pt]
e     & ::= & v \mid x          & \text{values, variables}\\
      & \mid & e_1 + e_2        & \text{addition (the whiteboard's }\oplus\text{)}\\
      & \mid & \textsf{let}\;x = e_1\;\textsf{in}\;e_2 & \text{binding}\\
      & \mid & c                & \text{a command (and not a control-flow)}\\
      & \mid & \textsf{if}\;e\;\textsf{then}\;e\;\textsf{else}\;e & \text{conditional}\\
\end{array}
\]



\subsection{Semantics :: of the Reagent Language}
\label{subsec:semantics}

Each configuration $C$ in $\Phi$ represents the global state of the program at a particular point during execution. Formally, a configuration is a six-tuple $C = \langle\sigma,\lambda,t,E,\hat{\sigma},\hat{E}\rangle$, where:

\begin{itemize}

\item \(\sigma : \Loc \rightarrow \mathbb{N}\) maps each shared memory location to a \(\mathbb{N}\).

\item \(\lambda \in \ThId \cup \{\bot\}\) records the \emph{thread} that currently holds the live transaction --- the lock owner. $\lambda = i$ means thread $i$ is inside its transaction and is the only thread permitted to take a transactional step; $\lambda = \bot$ means no transaction is active. Because at most one transaction runs at any time, $\lambda$ is a single global value, not a per-thread map.

\item \(t \in \TrId \cup \{\bot\}\) is the identifier of the live transaction --- a \emph{fresh} token minted by each $\bbegin$. It is what the emitted actions are tagged with ($\getev{l}{v}{t}$, $\eend\,@\,t$, and so on), and $t = \bot$ exactly when $\lambda = \bot$.

\item $E : \ThId \rightharpoonup \Expr$ maps each active thread to the
expression currently being evaluated by that thread.

\item $\hat{\sigma}: \Loc \rightarrow \mathbb{N}$ is the snapshot of the shared store taken when the live transaction began; it is defined only while $\lambda \neq \bot$, and is used to support transaction restart.

\item $\hat{E} \in \Expr$ is the expression from which execution should resume when a retry operation is encountered; like $\hat{\sigma}$ it is defined only while $\lambda \neq \bot$.

\end{itemize}

In every rule the \emph{active} thread is $i$ --- the one taking the step --- with
$E(i)=\ehole{r}$, where $r$ is the unique redex its evaluation context $\ectx$
locates (the subterm currently able to step). There is no scheduler: the
nondeterminism is inherent in $\rightarrow$, and any thread with an enabled redex
may be the one that steps. Every transactional rule
($\textsc{Get},\textsc{Set},\textsc{End},\textsc{Retry}$) carries the premise
$\lambda = i$: \emph{only the thread that owns the live transaction may step inside
it}. The conclusion rewrites that redex in place and \emph{frames} the rest of the
table, $E' = E[i \mapsto \ehole{\cdot}]$ with $E'(j) = E(j)$ for every $j \neq i$;
the label above the arrow is tagged with the live transaction id $t$ (the third
component of the configuration).
\begin{figure}[ht]
\small
\begin{mathpar}
\inferrule*[left=\textsc{Get}]{
  E(i)=\ehole{\getc{l}} \\ \lambda = i \\ \sigma(l)=v
}{
  \langle\sigma,\lambda,t,E,\hat{\sigma},\hat{E}\rangle
  \xrightarrow{\getev{l}{v}{t}}
  \langle\sigma,\lambda,t,E[i\mapsto\ehole{v}],\hat{\sigma},\hat{E}\rangle
}
\and
\inferrule*[left=\textsc{Set}]{
  E(i)=\ehole{\setc{l}{v}} \\ \lambda = i
}{
  \langle\sigma,\lambda,t,E,\hat{\sigma},\hat{E}\rangle
  \xrightarrow{\setev{l}{v}{t}}
  \langle\sigma[l\mapsto v],\lambda,t,E[i\mapsto\ehole{\unit}],\hat{\sigma},\hat{E}\rangle
}
\and
\inferrule*[left=\textsc{Begin}]{
  E(i)=\ehole{\bbegin\,;\,e}
  \\
  \lambda=\bot
  \\
  t\ \text{fresh}
}{
  \langle\sigma,\bot,\bot,E,\hat{\sigma},\hat{E}\rangle
  \xrightarrow{\bbegin\,@\,t}
  \langle
    \sigma,
    i,
    t,
    E[i\mapsto\ehole{\unit}],
    \sigma,
    e
  \rangle
}
\and
\inferrule*[left=\textsc{End}]{
  E(i)=\ehole{\eend} \\ \lambda = i
}{
  \langle\sigma,\lambda,t,E,\hat{\sigma},\hat{E}\rangle
  \xrightarrow{\eend\,@\,t}
  \langle\sigma,\bot,\bot,E[i\mapsto\ehole{\unit}],\hat{\sigma},\hat{E}\rangle
}
\and
\inferrule*[left=\textsc{Retry}]{
  E(i)=\ehole{\retry} \\ \lambda = i
}{
  \langle\sigma,\lambda,t,E,\hat{\sigma},\hat{E}\rangle
  \xrightarrow{\retry\,@\,t}
  \langle
    \hat{\sigma},
    \bot,
    \bot,
    E[i\mapsto\ehole{\bbegin\,;\,\hat{E}}],
    \hat{\sigma},
    \hat{E}
  \rangle
}
\end{mathpar}
\caption{Semantics of Reagent commands. Each rule fires on the redex of the
active thread $i$; the transactional rules require $\lambda = i$ (only the
owner thread may step), and the emitted label is tagged with the live transaction
id $t$. Unlisted six-tuple fields are unchanged.}
\label{fig:semantics}
\end{figure}
Figure~\ref{fig:semantics} gives a formal description of the transition
rules. $\textrm{\sc{Get}}$ reads the value stored at a shared location
$l$. If the active thread $i$ owns the live transaction ($\lambda = i$) and $\sigma(l)=v$, then $\getc{l}$ returns the value $v$; the emitted action $\getev{l}{v}{t}$ records both the value $v$ read and the id $t$ of the live transaction. The store and transaction state remain unchanged.

$\textrm{\sc{Set}}$ updates a shared location. When the owner thread $i$ (with $\lambda = i$) executes
$\setc{l}{v}$, the value stored at location $l$ is updated to $v$ and the
command returns the unit value $\unit$; the emitted action $\setev{l}{v}{t}$ records the written value $v$ and the transaction id $t$.

$\textrm{\sc{Begin}}$ starts a transaction. No transaction may currently be alive ($\lambda = \bot$); transactions are therefore mutually exclusive, with at most one live transaction across all threads at any time. $\textrm{\sc{Begin}}$ mints a \emph{fresh} transaction id $t$ (the new live id), records the executing thread as the owner by setting $\lambda := i$, and saves the current store and continuation in $\hat{\sigma}$ and $\hat{E}$ respectively.

$\textsc{End}$ commits the current transaction, resetting both the owner $\lambda$ and the live transaction id $t$ to $\bot$. The store is unchanged.


$\textsc{Retry}$ is retrying with \emph{lock-releasing}. Instead of keeping the thread inside the transaction and spinning, it (i) rolls the store back to the snapshot $\hat{\sigma}$, so the abandoned attempt has no effect; (ii) sets $\lambda=\bot$ and clears the live transaction id, \emph{releasing the single critical section} so that another thread may now $\bbegin$ and run; and (iii) re-arms $E(i)$'s expression as $\bbegin\,;\,\hat{E}$, so that when $i$ next steps it must re-acquire the lock through $\textsc{Begin}$ before re-running the saved body.

Because every transactional rule requires $\lambda = i$, a thread can only read, write, commit, or retry inside a transaction it opened itself. A $\getc{l}$, $\setc{l}{e}$, $\eend$, or $\retry$ reached with $\lambda \neq i$ --- in particular one with no matching enclosing $\bbegin$ --- has no rule to fire, so the run gets stuck rather than taking an ill-formed step. Well-formedness is thus enforced by the operational semantics itself, not assumed (Section~\ref{sec:examples}).

\subsection{Evaluation contexts}

An evaluation context $\ectx$ is an expression containing a single
hole $[\cdot]$ that marks exactly where the next reduction step occurs. We give
the evaluation context its own metavariable $\ectx$ precisely to keep it
separate from $E$, which represents a mapping from thread identifiers
($\mathit{Tid}$) to expressions ($\mathit{Expr}$); the two are distinct objects
and should not be conflated.
\[
\ectx \;::=\; [\cdot]
\;\mid\; \ectx\,;\,e
\;\mid\; \textsf{let}\;x = \ectx\;\textsf{in}\;e
\;\mid\; \ectx + e \;\mid\; v + \ectx
\;\mid\; \ectx == e \;\mid\; v == \ectx
\;\mid\; \textsf{if}\;\ectx\;\textsf{then}\;e\;\textsf{else}\;e
\;\mid\; \setc{l}{\ectx}.
\]
We write $\ehole{e}$ for the expression obtained by plugging $e$ into the hole. Thus
$E(t)=\ehole{\bbegin}$ asserts that the redex thread $t$ is about to execute is
$\bbegin$, sitting inside the surrounding continuation $\ectx$.

\subsection{The command rules in context form}

Every command rule now fires on its redex inside an arbitrary context $\ectx$,
preserving the continuation. The configuration is the six-tuple
$\langle\sigma,\lambda,t,E,\hat{\sigma},\hat{E}\rangle$ of
Section~\ref{subsec:semantics}.

\begin{itemize}
\item $E(i)=\ehole{\bbegin}$ and $E'=E[i\mapsto \ehole{\unit}]$: the $\bbegin$ reduces to
$\unit$ \emph{in place}, leaving the continuation $\ectx$.
\item $\hat{E}'=\ehole{\unit}$: the saved resume point is $\ehole{\unit}$,
the body \emph{after} $\bbegin$, which is precisely what $\textsc{Retry}$
reinstates.
\end{itemize}
\textsc{Retry} rolls the store back to
the snapshot $\hat{\sigma}$, releases the lock ($\lambda = \bot$), and
re-arms $E(i)$ as $\bbegin\,;\,\hat{E}$ so that the thread must re-acquire the
lock through $\textsc{Begin}$ before re-running.

\subsection{Administrative reductions}


\[
\begin{aligned}
\ehole{\textsf{let}\;x = v\;\textsf{in}\;e} &\;\reduces\; \ehole{\subst{e}{x}{v}} &
\ehole{\unit \,;\, e} &\;\reduces\; \ehole{e} \\
\ehole{n_1 + n_2} &\;\reduces\; \ehole{n} \quad (n = n_1 + n_2) &
\ehole{\textsf{if}\;\textsf{true}\;\textsf{then}\;e_1\;\textsf{else}\;e_2}
   &\;\reduces\; \ehole{e_1} \\
\end{aligned}
\]
Here the notation is substitution $\subst{e}{x}{v}$ in the \textsf{let}
rather than map update.




%% =====================================================================
\section{Initial configuration and termination}
\label{sec:init-final}

The transition system $\mathcal{S}=(\Phi,\rightarrow)$ has an \emph{infinite}
state space: a configuration $\langle\sigma,\lambda,t,E,\hat{\sigma},\hat{E}\rangle$
ranges over all stores, all thread maps, all snapshots, and all saved bodies, and
there are infinitely many of each. A program does not fix a single starting point;
it determines a \emph{set} of initial configurations, and a run may begin from any
one of them.

\subsection{Initial configurations}

The set of \emph{initial configurations} is
\[
\mathcal{I} \;=\;
\bigl\{\, \langle\, \sigma_0,\ \bot,\ \bot,\ E_0,\ \sigma_0,\ \unit \,\rangle
\;:\; \sigma_0, E_0 \text{ meet the conditions below} \,\bigr\}
\;\subseteq\; \Phi ,
\]
and a run may start from any $C_0 \in \mathcal{I}$.
\begin{itemize}
\item \emph{Locations are fixed and finite.} $\Loc$ is a fixed finite set, and
$\sigma_0 : \Loc \rightarrow \mathbb{N}$ is any initial store over it, so
$\mathcal{I}$ genuinely varies with $\sigma_0$.
\item \emph{Threads are fixed and finite.} $dom(E_0)$ is a fixed finite set of
thread identifiers; no thread is ever created or destroyed, so $dom(E) = dom(E_0)$
at every configuration of the run.
\item \emph{Every thread starts with a real program.} $E_0(i) \neq \unit$ for each
$i \in dom(E_0)$: a thread already at the unit value is finished and cannot step,
so an initial configuration must give every thread an unreduced expression for the
system to make progress.
\item \emph{No transaction is live.} $\lambda = \bot$ and $t = \bot$; accordingly
the snapshot and saved body start at $\hat{\sigma}_0 = \sigma_0$ and
$\hat{E}_0 = \unit$.
\end{itemize}

Fixing any $C_0 \in \mathcal{I}$, the configurations that matter are those
reachable from it,
\[
\mathrm{Reach}(C_0) \;=\; \{\, C \;:\; C_0 \rightarrow^{*} C \,\},
\]
where $\rightarrow^{*}$ is the reflexive--transitive closure of $\rightarrow$.
There is no scheduler: the nondeterminism of $\rightarrow$ is inherent, and at each
configuration \emph{any} thread whose redex is enabled may take the next step.
These reachable step sequences are the \emph{executions} of
Section~\ref{sec:fol-int}.


\subsection{Termination}

We do not define a ``final'' or ``terminal'' configuration; the transition
relation has no distinguished stopping states. We instead \emph{assume} that a run
halts after finitely many steps, and record the condition under which it has run
to completion: every thread has reduced to the unit value,
\[
E(i) = \unit \qquad \forall i \in dom(E).
\]

Here $\unit$ is the unit \emph{value} of the abstract syntax, so this is a
pointwise statement $E(i)=\unit$; writing $E : \ThId \rightharpoonup \{\unit\}$
would be a category error, since the grammar produces values, not sets.



%% =====================================================================
\section{Worked example: well-formedness is structural}
\label{sec:examples}

Because $\lambda$ now names the \emph{owner thread} and every transactional rule
carries the premise $\lambda = i$, a thread can take a transactional step only
inside a transaction it has itself opened. Call a program \emph{well-formed} when
every $\getc{l}$, $\setc{l}{e}$, $\eend$, or $\retry$ it reaches sits inside a
matching enclosing $\bbegin\dots\eend$. The point of this section is that
well-formedness is \emph{structural}: an ill-formed redex simply has no rule to
fire, so the run gets stuck rather than performing the action. We never need
well-formedness as a separate side condition.

\paragraph{Conventions.}
There is a single shared location $l$; we write $\st{n}$ for the store with
$\sigma(l)=n$. The configuration is the six-tuple
$\langle\sigma,\lambda,t,E,\hat{\sigma},\hat{E}\rangle$ of
Section~\ref{subsec:semantics}, with $\lambda$ the owner thread and $t$ the live
transaction id. Thread identifiers are $1,2\in\ThId$; $E$ is shown as a map
$\emap{1{:}e_1,\,2{:}e_2}$. Every run starts from
$C_0=\langle \st{0},\ \bot,\ \bot,\ E_0,\ \st{0},\ \unit\rangle$.

\subsection{An ill-formed \texttt{get} cannot fire}

Let thread $1$ run the ill-formed program $P \equiv \getc{l}$ --- a read with no
enclosing $\bbegin$. The initial configuration is
\[
C_0 \;=\; \langle \st{0},\ \bot,\ \bot,\ \emap{1{:}\getc{l}},\ \st{0},\ \unit\rangle .
\]
The active thread is $1$ and its redex is $\getc{l}$ (context $\ectx=[\cdot]$).
We check every rule of Figure~\ref{fig:semantics}:
\begin{itemize}
\item $\textsc{Get}$ requires $\lambda = 1$, but $\lambda = \bot$. \emph{Fails.}
\item $\textsc{Begin}$ requires the redex to be $\bbegin\,;\,e$, but it is
$\getc{l}$. \emph{Fails.}
\item $\textsc{Set},\textsc{End},\textsc{Retry}$ require their own redexes
($\setc{l}{v}$, $\eend$, $\retry$) and $\lambda = 1$; both premises fail.
\end{itemize}
No rule matches, so $C_0$ takes \emph{no} step:
\[
\langle \st{0},\ \bot,\ \bot,\ \emap{1{:}\getc{l}},\ \st{0},\ \unit\rangle
\ \not\rightarrow .
\]
The ill-formed $\getc{l}$ never executes; the thread is stuck. Contrast the
well-formed version, $\bbegin\,;\,(\getc{l}\,;\,\eend)$. There the leading
$\bbegin$ fires first (setting $\lambda:=1$), and only \emph{then}, with
$\lambda = 1$, can the $\getc{l}$ fire: it is exactly the $\bbegin$ that grants
the ownership $\getc{l}$ demands.

\subsection{A thread cannot act inside another thread's transaction}

This is the case the premise $\lambda = i$ is really designed to rule out. Take
two threads,
\[
E_0 \;=\; \emap{1{:}\bbegin\,;\,(\getc{l}\,;\,\eend),\ 2{:}\getc{l}} ,
\]
so thread $1$ is well-formed and thread $2$ tries a bare $\getc{l}$. Let thread
$1$ open its transaction:
{\small
\[
\begin{aligned}
&\langle \st{0},\ \bot,\ \bot,\ E_0,\ \st{0},\ \unit\rangle\\
&\quad\xrightarrow{\ \bbegin\,@\,t_1\ }\
\langle \st{0},\ 1,\ t_1,\ \emap{1{:}\unit\,;\,(\getc{l}\,;\,\eend),\ 2{:}\getc{l}},\ \st{0},\ \getc{l}\,;\,\eend\rangle .
\end{aligned}
\]}%
Now $\lambda = 1$: thread $1$ owns the live transaction $t_1$. Suppose thread $2$
is the next to step; its redex is $\getc{l}$. The only rule that could
fire is $\textsc{Get}$, and it requires $\lambda = 2$ --- but $\lambda = 1$. The
premise fails, so thread $2$'s read is \emph{not enabled}:
\[
\text{there is no } C' \text{ with }\quad
\langle \st{0},\ 1,\ t_1,\ \emap{1{:}\dots,\ 2{:}\getc{l}},\ \st{0},\ \getc{l}\,;\,\eend\rangle
\ \xrightarrow{\ \getev{l}{0}{\,\cdot\,}\ }\ C' .
\]
Under the old rule (premise $\lambda = t$, with $t$ unconstrained by the acting
thread) this step \emph{would} have fired, letting thread $2$ read inside thread
$1$'s transaction and mislabel it $@t_1$. The premise $\lambda = i$ removes that:
only the owner may step, so thread $2$ must wait until thread $1$ commits
($\lambda$ returns to $\bot$) and then open its own transaction. Well-formedness
holds by construction.

%% =====================================================================
\section{Internal consistency}
\label{sec:fol-int}

We now package the operational behaviour of Section~\ref{subsec:semantics} into
the declarative vocabulary of Section~\ref{sec:declarative}, so that the internal
consistency axiom $\textsc{Int}$ can be read off it. Throughout, $C$ ranges over
configurations $\langle\sigma,\lambda,t,E,\hat{\sigma},\hat{E}\rangle$ and
$\rightarrow$ is the labelled transition of Figure~\ref{fig:semantics}; our notion
of \emph{history} follows Herlihy and Wing~\cite{HerlihyWing1990}.

\subsection{Executions, traces, and histories}

\begin{definition}[Execution]
An \emph{execution} from an initial configuration $C_0 \in \mathcal{I}$ is a finite
sequence of \emph{steps}
\[
C_0 \longrightarrow C_1 \longrightarrow \cdots \longrightarrow C_n ,
\]
in which every step $C_k \longrightarrow C_{k+1}$ is derivable from
Figure~\ref{fig:semantics} and is either a \emph{transition} (a command rule
$\textsc{Get},\textsc{Set},\textsc{Begin},\textsc{End},\textsc{Retry}$) or an
\emph{administrative reduction} (Section~\ref{subsec:semantics}).
\end{definition}

\begin{definition}[Trace]
The \emph{trace} of an execution is the sequence of labels emitted along it,
\[
\tau \;=\; \ell_1 \cdot \ell_2 \cdots \ell_n ,
\]
where a transition step contributes its \emph{visible action} --- one of
$\bbegin\,@\,t$, $\eend\,@\,t$, $\retry\,@\,t$, $\getev{l}{v}{t}$, $\setev{l}{v}{t}$
--- and an administrative step contributes the silent label $\tau_{\mathrm{adm}}$.
The trace thus records \emph{where} administrative reductions occur, marked by
$\tau_{\mathrm{adm}}$.
\end{definition}

\begin{definition}[History]
Following Herlihy and Wing~\cite{HerlihyWing1990}, the \emph{history} $H(\tau)$ of a
trace $\tau$ is its subsequence of transaction events: each transaction's
\emph{invocation} $\bbegin\,@\,t$ and, when it commits, its matching
\emph{response} $\eend\,@\,t$ (an aborted attempt is marked by $\retry\,@\,t$). The
intervening $\getev{l}{v}{t}$ and $\setev{l}{v}{t}$ actions are the transactions'
internal effects.
\end{definition}



\subsection{Transactions and commitment}

\begin{definition}[operation sequence]
The \emph{operation sequence} $\elided(\tau)$ of a trace $\tau$ is $\tau$ with
every silent $\tau_{\mathrm{adm}}$ deleted --- the visible actions, in order:
\[
\elided(\tau) \;=\; \ell_{k_1} \cdot \ell_{k_2} \cdots \ell_{k_m},
\qquad \ell_{k_j} \neq \tau_{\mathrm{adm}} .
\]
For instance,
\[
\elided(\tau) \;=\;
\bbegin\,@\,t_1 \cdot \getev{l}{0}{t_1} \cdot \setev{l}{1}{t_1}
\cdot \eend\,@\,t_1 \cdot \bbegin\,@\,t_2 \cdots
\]
\end{definition}

For a transaction identifier $t$ we write $\elided(\tau)|_t$ for the subsequence of
$\elided(\tau)$ consisting of exactly the actions tagged $@\,t$, in order. Two
structural facts hold. First, $\textsc{Begin}$ allocates a \emph{fresh} identifier
at each opening (Figure~\ref{fig:semantics}), so a given $t$ is opened at most once:
$\bbegin\,@\,t$ occurs at most once in $\elided(\tau)$. Second, because at most one
transaction holds the lock at a time, the actions of
$t$ form one \emph{contiguous} block of $\elided(\tau)$, opened by that
$\bbegin\,@\,t$ and closed by a single terminator that is either $\eend\,@\,t$ or
$\retry\,@\,t$:
\[
\elided(\tau)|_t \;=\; \bbegin\,@\,t \,\cdot\, body \,\cdot\, c_t,
body \in \bigl\{\,\getev{l}{v}{t},\ \setev{l}{v}{t} \,\bigr\}^{*},
c_t \in \bigl\{\, \eend\,@\,t,\ \retry\,@\,t \,\bigr\}.
\]

\begin{definition}[Committed transaction]
\label{def:committed-tx}
A transaction identifier $t$ is \emph{committed} in $\tau$ when its opening
$\bbegin\,@\,t$ is matched by a closing end, i.e.\ its terminator is
\[
c_t \;=\; \eend\,@\,t.
\]
If instead $c_t = \retry\,@\,t$ the transaction is \emph{uncommitted}; if $t$ is opened
but the trace stops before any terminator (no $c_t$ occurs) it is \emph{pending}.
\end{definition}

A single logical transaction may take several \emph{attempts}: each $\retry\,@\,t$
aborts an attempt and re-arms the \emph{same} thread with $\bbegin\,;\,\hat{E}$ to
open a fresh one (Figure~\ref{fig:semantics}). So the attempts of one logical
transaction form a chain of aborted tries ending in a committing one; for example
\[
\bbegin\,@\,t_1 \,\cdot\, \retry\,@\,t_1 \,\cdot\,
\bbegin\,@\,t_2 \,\cdot\, \retry\,@\,t_2 \,\cdot\,
\bbegin\,@\,t_3 \,\cdot\, \eend\,@\,t_3
\]
is \emph{one} completed transaction: $t_1,t_2$ are discarded attempts and the
commit at $t_3$ is precisely the result of retrying them.

To make this precise, let $\mathrm{thr}(t)$ be the thread that executed
$\bbegin\,@\,t$. Because $\textsc{Retry}$ re-arms that same thread, an aborted
attempt has a unique \emph{retry-successor}:
\[
t \rightsquigarrow t'
\quad\text{iff}\quad
c_t = \retry\,@\,t,\ \ \mathrm{thr}(t') = \mathrm{thr}(t),
\ \text{and } t' \text{ is the next transaction } \mathrm{thr}(t) \text{ opens.}
\]
A committed attempt has no successor (an $\eend$ stops the thread's chain), so
$\rightsquigarrow$ is a partial function and its maximal chains
\[
t_1 \rightsquigarrow t_2 \rightsquigarrow \cdots \rightsquigarrow t_k
\]
partition the attempts of $\tau$. Each maximal chain is one \emph{logical
transaction}; since only its last attempt can carry an $\eend$, the chain is
\emph{completed} exactly when $c_{t_k} = \eend\,@\,t_k$ (and then $t_1,\dots,t_{k-1}$
are all aborted attempts).

\begin{definition}[Committed history]
In the sense of Herlihy and Wing~\cite{HerlihyWing1990}, a history is
\emph{complete} when every invocation has a matching response. Here that is exactly
the condition that every logical transaction be \emph{completed}: every maximal
$\rightsquigarrow$-chain ends in a committing attempt $\eend\,@\,t_k$ --- no attempt
is left \emph{pending} (the execution does not stop inside an attempt) and no
aborted attempt is left \emph{dangling} (every $\retry\,@\,t$ has a successor). The
\emph{committed history} is the set of these complete histories,
\[
\mathcal{H}_{\mathrm{committed}} \;=\;
\bigl\{\, H(\tau) \;:\; \text{$\tau$ a trace whose every logical transaction is completed} \,\bigr\}.
\]
This \emph{permits} $\retry\,@\,t$ actions: they are aborted attempts on the way to
a commit, not a violation.
\end{definition}

\begin{remark}
Permitting $\retry\,@\,t$ is not the same as tolerating an \emph{unfinished} retry
chain. Suppose one thread commits while another is still retrying:
\[
\elided(\tau) \;=\;
\bbegin\,@\,t \,\cdot\, \eend\,@\,t \,\cdot\,
\bbegin\,@\,u_1 \,\cdot\, \retry\,@\,u_1 \,\cdot\,
\bbegin\,@\,u_2 \,\cdot\, \retry\,@\,u_2 .
\]
The chain of $t$ is completed, but the chain $u_1 \rightsquigarrow u_2$ ends in a
$\retry$ with no successor, so it is \emph{not} completed. Since a committed history
requires \emph{every} chain to complete, we have
$H(\tau) \notin \mathcal{H}_{\mathrm{committed}}$; the commit of $t$ does not
rescue it. Only once the $u$-thread re-runs and reaches an $\eend$, settling its
chain, does the extended trace become committed.
\end{remark}


\subsection{Lifting Cerone's \textsc{Int}}

On the committed history the internal-consistency axiom $\textsc{Int}$ of Cerone et
al.~\cite{CeroneBernardiGotsman2015} (Section~\ref{sec:declarative}) holds
verbatim; the pessimistic single lock makes the lift immediate. Two facts do the
work.

\begin{itemize}
\item \emph{Isolation.} Inside a committed transaction $t$, run by its owner
thread $i$ between $\bbegin\,@\,t$ and $\eend\,@\,t$, the premise $\lambda=i$
shared by every transactional rule means no other thread can step. Hence the value each $\getev{l}{v}{t}$ observes is
exactly the one written or read by the $\po$-latest earlier action on $l$, which
is precisely the conclusion of $\textsc{Int}$.
\item \emph{Retry is effect-free.} If an attempt aborts, $\textsc{Retry}$ rolls
the store back to the snapshot $\hat{\sigma}$ and releases the lock, so the aborted
events leave no mark on the committed store; and on re-entry the first
$\getev{l}{\cdot}{t'}$ of the fresh attempt re-reads the current committed value.
Thus the reads a committed transaction ultimately performs are the consistent
ones, and the earlier $\getev{l}{v}{t}$ of an aborted attempt stays consistent for
its own (discarded) attempt.
\end{itemize}

Consequently, dropping the aborted attempts and passing to
$\mathcal{H}_{\mathrm{committed}}$ loses no committed behaviour, and
$\textsc{Int}$ transfers directly to the Reagent runs: for every committed
transaction and every $\getev{l}{v}{t}$ in it, $v$ is the value of the
$\po$-maximal earlier read-or-write on $l$.


%% =====================================================================
\section{A declarative semantics for internal consistency}
\label{sec:declarative}






%%
WIP
%%
(still trying to figure it out)
=======
Fix a single run as a history: a set of events $E$ totally ordered by program
order $\po$, with the retry events written $R(E) \subseteq E$. Because at most
one transaction is live under mutual exclusion; the lock is held across
retries;  one such $E$ suffices, so we bind it once here and let every
quantifier below range over it.

\[
\forall e \in E.\ \forall l,v.\quad
\bigl(e = (\_,\rdev{l}{v}) \,\wedge\, \po^{-1}(e)\cap\HEvent_l \neq \emptyset\bigr)
\implies
\maxpo\bigl(\po^{-1}(e)\cap\HEvent_l\bigr) = (\_,\anyop{l}{v}),
\tag{\textsc{Int}}
\]

where $\anyop{l}{v}$ abbreviates ``a read \emph{or} a write of $v$ on $l$''.

This has been modified to the following specification.

\[
\att_l(e) \;=\; \bigl\{\, f \;\mid\; f \pob e \;\wedge\; f \in \HEvent_l \;\wedge\;
\neg\,\exists\, r \in R(E).\ f \pob r \pob e \,\bigr\}
\]
\[
\forall e \in E.\ \forall l,v.\quad
\bigl(e = (\_,\rdev{l}{v}) \,\wedge\, \att_l(e) \neq \emptyset\bigr)
\implies
\maxpo\bigl(\att_l(e)\bigr) = (\_,\anyop{l}{v}).
\tag{\textsc{Int}$^{r}$}
\]

Dropping visibility and letting a read observe the arbitration-latest earlier
writer gives the external-consistency axiom in its strict-serializable form
(\textsc{Vis} is total and coincides with \textsc{Ar}, so only \textsc{Ar}
remains):
\[
\begin{aligned}
\forall T \in H.\ \forall x,n.\quad T \vdash \mathsf{Read}\ x : n \implies {}
&\bigl(\{\, S \in H \mid S <_{\mathrm{ar}} T \,\wedge\, S \vdash \mathsf{Write}\ x : \_ \,\} = \emptyset \,\wedge\, n = 0\bigr)\\[2pt]
&{}\vee\ \mathrm{max}_{\mathrm{ar}}\{\, S \in H \mid S <_{\mathrm{ar}} T \,\wedge\, S \vdash \mathsf{Write}\ x : \_ \,\} \vdash \mathsf{Write}\ x : n,
\end{aligned}
\tag{\textsc{Ext}}
\]



WIP



In my understanding the abstract syntax doesn't change, because \texttt{get(l,v)@t} and \texttt{set(l,v)@t} are runtime actions. The value $v$ is the value
observed at execution time and t is the active transaction identifier. Neither of which a programmer can write, so the source commands stay \texttt{get(l)} and \texttt{set(l,v)}.


\begin{thebibliography}{9}

\bibitem{HerlihyWing1990}
Maurice P.\ Herlihy and Jeannette M.\ Wing.
\newblock Linearizability: a correctness condition for concurrent objects.
\newblock \emph{ACM Trans.\ Program.\ Lang.\ Syst.}, 12(3):463--492, 1990.

\bibitem{CeroneBernardiGotsman2015}
Andrea Cerone, Giovanni Bernardi, and Alexey Gotsman.
\newblock A framework for transactional consistency models with atomic visibility.
\newblock In \emph{CONCUR 2015}, LIPIcs 42, pages 58--71, 2015.

\end{thebibliography}

\end{document}
